\documentclass[a4paper,10pt]{article}
\newcommand\subdir{/home/koen/LaTeX-setup/}
\input{\subdir/LaTeX-files/imports.tex}
\input{\subdir/LaTeX-files/master-internship.tex}
\usepackage{\subdir/LaTeX-files/aastex}
\usepackage{natbib}
\bibliographystyle{\subdir/LaTeX-files/astroadstitle.bst}

%Set the title, Author and date
\title{Notes Week 36}
\author{Koen Essers}
\tutorial{Master Internship}
\date{\today}

%Set the header style
\header{\tutorialstring}

%Begin the document
\begin{document}
\setuplayout
\section{Dense Mass Spacing Results}

Below I plot the envelope mass for the new models and the old ones.

\begin{figure}[ht]
	\centering
	\incplot{w36-mass-at-TP}
\end{figure}

Quite clearly, even for \(0.01\; M_\odot \) differences in mass, there is no predictable relation between
changes in initial mass and envelope mass during the last thermal pulse.

It seems to be very ``chaotic''.
Small perturbations in initial mass lead to quite drastically different
final TPAGB evolutions.

To highlight this, let me show the last few thermal pulses for the different masses between
\(1.9\; M_\odot \) and \(2.1\; M_\odot \).

\begin{figure}[ht]
	\centering
	\incplot{w36-diff-radial-evolution}
\end{figure}

\section{Moe and Di Stefano}

I use a standard \citet{salpeter1955} IMF.
Differences between this IMF and for example, \citet{chabrier1903} should be very minor, since we are only interested in primaries that produce
enough s-process elements in third dredge-up episodes.
This happens at masses of about \(1.2\; M_\odot \) and higher.
We also only go up to masses of \(M = 3.0\; M_\odot \).
Above this mass, it is difficult to reproduce barium dwarfs with
masses within the observed range.
Therefore, I have not simulated these TPAGB stars
with MESA for masses above \(3.0\; M_\odot \).

For the plots below, I only show periods below \(30\,000\) days, because periods above this
limit will not produce barium dwarfs within the observed period range, and \emph{definitely}
not RLOF TPAGB binaries.

\begin{figure}[ht]
	\marginfignote{0}{Here, I show a population of 153714(000) \emph{binaries} using the parameterization of
		\citet{moe2016}.
		Their paper explicitly describes multiplicities, and not only binaries,
		so really I generated a bunch of singles, binaries, triples and quadruples, but I only show
		the binaries here.
	}
	\marginfignote{12}{Clearly, the mass distribution is similar to the \citet{salpeter1955}
		IMF.
		However, it is slightly flatter.
		There are slightly more heavier stars and
		slightly fewer lighter stars.
	}
	\marginfignote{20}{The binaries also contain many "twin" binaries.
		This is about 10 to 15\% of the total systems.
	}
	\marginfignote{25}{
		Lastly, we see that, in log space, there are many more wide orbits.
		However, if you would look at the systems in linear space, the peak is towards lower periods.
	}
	\centering
	\incplot{w36-population-1}
\end{figure}

\begin{figure}[ht]
	\marginfignote{0}{Here, I show the same distributions, but now for the population containing 91111(000) triples.}
	\marginfignote{5}{There are fewer triples, which is to be expected within this mass range.
		Let me be fully clear on what I do with the triples as well as the quadruples.
		Triples contain two binaries.
		I simply include none, one, or both of these orbits, only dependent on if they are within the period range.
		This means the fraction of triple \emph{orbits} over binary \emph{orbits} will be \emph{larger} than the fraction of triple \emph{systems} over binary \emph{systems}.
	}
	\marginfignote{23}{The mass-ratio distribution and period distributions are not very different.
		However, we see a much shallower IMF.
	}
	\centering
	\incplot{w36-population-trips}
\end{figure}

\begin{figure}[ht]
	\marginfignote{0}{Here, I show the same distributions, but now for the population containing 27730(000) triples.}
	\marginfignote{5}{The mass-ratio distribution and period distributions are not very different.
		However, we see a much shallower IMF, even shallower than for the triples.
	}
	\centering
	\incplot{w36-population-quads}
\end{figure}

\nextpage

\begin{figure}[ht]
	\centering
	\incplot{w36-compare-distribution-sal}
\end{figure}

\begin{figure}[ht]
	\centering
	\incplot{w36-compare-distribution-bhl}
\end{figure}

\begin{figure}[ht]
	\centering
	\incplot{w36-compare-distribution-fast}
\end{figure}

\begin{figure}[ht]
	\marginfignote{0}{If we require the binaries to have \(q > 0.5\) for stable mass transfer,
		we get slightly different distributions.
		Here I also increased the barium threshold to [s/Fe] > 0.5.
	}
	\centering
	\incplot{w36-compare-distribution-sal-high-q}
\end{figure}

\begin{figure}[ht]
	\centering
	\incplot{w36-compare-distribution-bhl-high-q}
\end{figure}

\begin{figure}[ht]
	\centering
	\incplot{w36-compare-distribution-fast-high-q}
\end{figure}

\begin{figure}[ht]
	\marginfignote{0}{\(\varepsilon =0.1\)}
	\centering
	\incplot{w36-compare-distribution-sal-high-q-correct-mass}
\end{figure}

\nextpage

\begin{figure}[ht]
	\marginfignote{0}{\(\varepsilon =0.0\)}
	\marginfignote{2}{The fast winds models cannot produce barium stars.
		This is simply because
		we \emph{only} have wind mass transfer, and the fast winds do not transfer mass.
	}
	\marginfignote{10}{For the Saladino winds, we also see that the lower initial periods are now
		missing.
		This is because these models do not experience wind before RLOF occurs.
		In the simple case
		explored here, mass is stripped in its entirety during RLOF.
		However, we have seen that this is
		not very realistic.
		Therefore, the plot above is probably much more relevant.
	}
	\centering
	\incplot{w36-compare-distribution-sal-high-q-correct-mass-eps0}
\end{figure}

\begin{figure}[ht]
	\centering
	\incplot{w36-compare-distribution-bins-trips-quads}
\end{figure}

\begin{figure}[ht]
	\centering
	\incplot{w36-s-fe}
\end{figure}

\begin{figure}[ht]
	\centering
	\incplot{w36-m_env}
\end{figure}

\begin{figure}[ht]
	\centering
	\incplot{w36-s-m_env}
\end{figure}

\begin{figure}[ht]
	\centering
	\incplot{w36-s-m_env-bhl}
\end{figure}

\begin{figure}[ht]
	\centering
	\incplot{w36-s-m_env-f}
\end{figure}

\nextpage
\begin{figure}[ht]
	\centering
	\incplot{w36-s-m}
\end{figure}

\begin{figure}[ht]
	\centering
	\incplot{w36-s-p}
\end{figure}

\begin{figure}[ht]
	\centering
	\incplot{w36-m-m_env-p}
\end{figure}

\begin{figure}[ht]
	\centering
	\incplot{w36-m-m_env-winds}
\end{figure}

\bibliography{master.bib}
\nextpage
\appendix

\section{Moe and Di Stefano Tests}

\begin{figure}[ht]
	\marginfignote{0}{Compared with Figure 35 from \citet{moe2016}, these results look sensible.}
	\centering
	\incplot{w36-test-F_twin}
\end{figure}

\begin{figure}[ht]
	\marginfignote{0}{Compared with Figure 35 from \citet{moe2016}, these results look sensible.}
	\centering
	\incplot{w36-test-gamma-large}
\end{figure}

\begin{figure}[ht]
	\marginfignote{0}{Compared with Figure 35 from \citet{moe2016}, these results look sensible.}
	\centering
	\incplot{w36-test-gamma-small}
\end{figure}

\begin{figure}[ht]
	\marginfignote{0}{Compared with Figure 2 from \citet{moe2016}, these results look sensible.}
	\centering
	\incplot{w36-test-q-sample}
\end{figure}

\begin{figure}[ht]
	\marginfignote{0}{Compared with Figure 36 from \citet{moe2016}, these results look sensible.}
	\centering
	\incplot{w36-test-eta}
\end{figure}

\begin{figure}[ht]
	\marginfignote{0}{Based on how power laws should look and the figure above, I think this makes sense.}
	\centering
	\incplot{w36-test-ecc-sample}
\end{figure}

\begin{figure}[ht]
	\marginfignote{0}{Compared with Figure 37 from \citet{moe2016}, these results look sensible.}
	\centering
	\incplot{w36-test-log_P_q_gt_0_3}
\end{figure}

\begin{figure}[ht]
	\marginfignote{0}{Compared with Figure 37 from \citet{moe2016}, these results look sensible.}
	\centering
	\incplot{w36-test-log_P}
\end{figure}

\begin{figure}[ht]
	\marginfignote{0}{Sampling reproduces the period frequencies well.}
	\centering
	\incplot{w36-test-period-sampling}
\end{figure}

\begin{figure}[ht]
	\marginfignote{0}{Compared with Figure 39 from \citet{moe2016}, these results look sensible.}
	\centering
	\incplot{w36-multiplicity}
\end{figure}

\nextpage

\section{Additional Plots}

\begin{figure}[ht]
	\marginfignote{0}{Binaries, \(\varepsilon =0.1\), BHL}
	\centering
	\incplot{w36-m_env_bhl}
\end{figure}

\begin{figure}[ht]
	\marginfignote{0}{Binaries, \(\varepsilon =0.1\), fast}
	\centering
	\incplot{w36-m_env_f}
\end{figure}

\end{document}
